% TOP_PROBLEM: {"id": 30006811, "title": "Banach-Mazur Rotation Problem", "category_id": 8, "category": {"id": 8, "name": "analysis", "display_name": "Analysis", "description": "Limits, continuity, calculus, and function theory.", "slug": "analysis", "order_index": 8, "created_at": "2026-07-31T15:26:25.671Z"}, "status": "open"}
% FIRST_POSTED: 2026-09-27
% !TeX program = lualatex
% ATTEMPT_STATUS: unresolved
\documentclass[11pt]{article}
\usepackage[margin=25mm]{geometry}
\usepackage{fontspec}
\setmainfont{Latin Modern Roman}
\usepackage{amsmath,amssymb,mathtools}
\usepackage{unicode-math}
\setmathfont{Latin Modern Math}
\usepackage{xurl,hyperref}
\hypersetup{colorlinks=true,urlcolor=blue}
\setcounter{secnumdepth}{0}
\setlength{\parindent}{0pt}
\setlength{\parskip}{5pt}
\setlength{\emergencystretch}{3em}
\begin{document}
\section*{\texorpdfstring{Banach-Mazur Rotation Problem}{Banach-Mazur Rotation Problem}}\label{30006811}
\subsection{Definitions and mathematical statement}
A real Banach space is a complete normed real vector space. Its surjective linear isometries are linear bijections $U:X\to X$ with $\|Ux\|=\|x\|$. Transitivity on the unit sphere means
\[
 \forall x,y\in X,\quad\|x\|=\|y\|=1
 \quad\Longrightarrow\quad\exists U\text{ as above with }Ux=y.
\]
For infinite-dimensional separable $X$, does this force the norm to arise from an inner product, equivalently the parallelogram identity
$\|x+y\|^2+\|x-y\|^2=2\|x\|^2+2\|y\|^2$?
The conclusion is linear \emph{isometry} to a Hilbert space, not merely a bounded linear isomorphism. Approximate transitivity is not the selected hypothesis.
\par\smallskip\textit{Formulation and concept references:} \href{https://arxiv.org/abs/2012.08344}{[S1]}.

\subsection{Short English statement}
If every unit vector in a separable real Banach space can be carried exactly to every other by a linear isometry, must its norm be Hilbertian?
\subsection{Research attempt}
\paragraph{Scope and source check.}
This attempt was developed by GPT-6 Astra Ultra on 27 September 2026.
The exact, real, separable, infinite-dimensional question is retained.
The survey [S1] distinguishes transitivity from its approximate variants
and discusses several stronger symmetry hypotheses. None of those
stronger hypotheses is silently imposed here. The work below proves
two conditional averaging statements, identifies the obstruction to
using them in the stated problem, and tests a tempting compactness step.
These are partial arguments, not a claim of a new resolution.

Write $G$ for the group of all surjective linear isometries of $X$.
The desired conclusion is an identity for the given norm, not the
existence of some other equivalent Hilbert norm. A possible proof route
is to construct a positive definite $G$-invariant quadratic form and
then use sphere transitivity. A possible disproof route would require
a transitive norm with a nonzero parallelogram defect. Merely finding
a highly symmetric non-Hilbertian norm does not meet that requirement.

\paragraph{An invariant quadratic form would settle the question.}
Suppose there is a real symmetric bilinear form $B$ on $X$ satisfying
\[
 a\|x\|^2\le B(x,x)\le b\|x\|^2,
 \qquad B(Ux,Uy)=B(x,y)\quad(U\in G),                 \tag{1}
\]
where $0<a\le b<\infty$. Fix a norm-unit vector $e$.
Transitivity gives, for each norm-unit $x$, an isometry $U$ with
$Ue=x$, hence $B(x,x)=B(e,e)=c>0$. Homogeneity now gives
\[
                         B(x,x)=c\|x\|^2             \tag{2}
\]
for every $x$, including zero. Bilinearity of $B$ implies the
parallelogram law, and $B/c$ is the required inner product.
Completeness follows from completeness of the original norm.
Thus constructing (1) is sufficient; it is not already contained in
the transitivity assumption.

The same conclusion holds if the action is only approximately
transitive and $B$ is continuous: approximate an arbitrary unit vector
by an orbit and pass to the limit in $B(x,x)$. This observation does not
say that approximate transitivity by itself produces such a form.

\paragraph{The finite-dimensional case, with the averaging justified.}
Let temporarily $\dim X=d<\infty$ and choose any Euclidean inner product
$\langle\cdot,\cdot\rangle_0$. Its norm is equivalent to $\|\cdot\|$,
so all elements of $G$ and their inverses have uniformly bounded
Euclidean operator norm. The set $G$ is closed in the space of matrices:
if $U_j\to U$, then $\|Ux\|=\lim_j\|U_jx\|=\|x\|$, so $U$ is
injective and, in finite dimension, surjective. Consequently $G$ is a
compact group.

Use normalized Haar measure $\mu$ on this compact group and set
\[
              B(x,y)=\int_G\langle Ux,Uy\rangle_0\,d\mu(U).
                                                               \tag{3}
\]
If $a\|z\|^2\le|z|_0^2\le b\|z\|^2$, the integrand at $x=y$
lies between $a\|x\|^2$ and $b\|x\|^2$, uniformly in $U$.
Thus $B$ is positive definite with precisely the bounds in (1).
Right invariance of Haar measure proves $G$-invariance. Formula (2)
finishes the finite-dimensional proof.

The compact-group integration theorem is the external standard input
in this argument. Compactness is checked above; it cannot be asserted
for the isometry group in the infinite-dimensional problem.
In particular, choosing an arbitrary Euclidean structure separately
on each finite-dimensional subspace does not make (3) available:
an isometry carrying one vector to another need not preserve that
subspace, and the resulting quadratic forms need not be compatible
on intersections.

\paragraph{An infinite-dimensional conditional theorem.}
Assume, in addition to sphere transitivity, both of the following:
there is an equivalent Hilbert inner product $\langle\cdot,\cdot\rangle_0$
with $a\|x\|^2\le|x|_0^2\le b\|x\|^2$, and the abstract group $G$
has a right-invariant mean on all bounded real functions on $G$.
The latter means a positive linear functional $m$ with $m(1)=1$ and
\[
                m(f(\,\cdot\,V))=m(f)\quad(V\in G).
\]
This is an extra amenability assumption on the discrete group;
no equivalence with other topological amenability conventions is used.

For $x,y\in X$, the function
$f_{x,y}(U)=\langle Ux,Uy\rangle_0$ is bounded, since
$|f_{x,y}(U)|\le b\|x\|\|y\|$. Define $B(x,y)=m(f_{x,y})$.
Linearity of the mean gives bilinearity and symmetry; positivity gives
the two diagonal bounds. Also
\[
 f_{Vx,Vy}(U)=f_{x,y}(UV),\qquad B(Vx,Vy)=B(x,y).
\]
The preceding invariant-form lemma proves that the original norm is
Hilbertian. This is a complete conditional implication with explicit
additional assumptions. Neither the equivalent Hilbert structure nor
the invariant mean has been derived from separability and transitivity.

One cannot replace $f_{x,y}$ by the polarization expression of an
arbitrary norm and declare the average bilinear. Polarization is
bilinear exactly when the underlying squared norm satisfies the
parallelogram law, which is the conclusion being sought. Averaging
nonquadratic functions does not repair that missing premise without
an additional argument.

\paragraph{Why compactness cannot be recovered from bounded orbits.}
Put the strong operator topology on $G$, so $U_j\to U$ means convergence
of $U_jx$ in norm for each fixed $x$. If $G$ were compact in this topology,
the continuous orbit map $U\mapsto Ue$ would make the entire unit sphere
compact, by transitivity. This is impossible in infinite dimension:
Riesz's lemma constructs unit vectors $x_j$ with
\[
 \operatorname{dist}(x_j,\operatorname{span}\{x_1,\ldots,x_{j-1}\})
       >1/2,
\]
so no subsequence is Cauchy. Thus the most direct extension of the
finite-dimensional compact-group proof is ruled out by the hypothesis
that the dimension is infinite, even for Hilbert spaces.

A second failure concerns surjectivity under limits. In real
$\ell_2(\mathbb N)$ define $U_m$ by cycling the first $m$ basis vectors,
\[
 U_me_j=e_{j+1}\ (1\le j<m),\qquad U_me_m=e_1,
 \qquad U_me_j=e_j\ (j>m).
\]
Every $U_m$ is an orthogonal surjection. On finitely supported vectors
$U_m$ eventually agrees with the unilateral shift $Se_j=e_{j+1}$.
Uniform operator norm one and approximation by finitely supported
vectors give $U_mx\to Sx$ for every $x\in\ell_2$.
Yet $e_1$ is not in the range of $S$. A pointwise diagonal extraction
on a countable dense set, even if available, can therefore lose the
surjectivity required by the problem. Separability does not fix this.

\paragraph{Pair geometry is the remaining issue.}
For arbitrary $x,y$ put
\[
 \Delta(x,y)=\|x+y\|^2+\|x-y\|^2-2\|x\|^2-2\|y\|^2.
\]
One always has $\Delta(Ux,Uy)=\Delta(x,y)$.
Transitivity can normalize $x$ to a fixed unit vector, but it applies
the same $U$ to $y$. It does not make the stabilizer of $x$ transitive
on all possible second vectors or fix their relative geometry.
Consequently an argument which independently rotates both vectors
has strengthened the symmetry assumption.

For a concrete test, in $\ell_p^2$ the coordinate unit vectors give
\[
                 \Delta(e_1,e_2)=2^{1+2/p}-4.
\]
This vanishes precisely at $p=2$. Signed coordinate permutations
still act by isometries for every $p$, so a substantial visible
symmetry group does not imply the parallelogram law. This example
does not have the required sphere transitivity and is not a
counterexample. It tests only the proposed inference from a few
explicit rotations to full Hilbert geometry.

The same distinction appears in an orbit argument: checking images of
a dense set of unit vectors is not enough unless one constructs exact
surjective isometries for every target, with a limit mechanism that
preserves invertibility. The shift example shows the precise danger.

\paragraph{Verification and next bottleneck.}
The finite-dimensional proof includes closedness, compactness,
positivity, and invariance, rather than just an instruction to average.
The conditional infinite-dimensional proof uses a mean only on bounded
functions and checks its right-translation convention. The shift
countertest is verified on the basis and extended by a uniform bound.
The scalar defect formula supplies an independent elementary check.
No numerical computation of finitely many directions can establish
transitivity or the parallelogram law on an arbitrary Banach space.

\paragraph{Final status.}
\textbf{UNRESOLVED.} The strongest results established here are the
finite-dimensional assertion and the invariant-mean implication under
the two stated extra assumptions. The first missing step is a
construction, from the actual hypotheses alone, of a bounded positive
definite invariant bilinear form as in (1), or another argument forcing
$\Delta(x,y)=0$ for all pairs. Concrete next targets are compatible
quadratic forms on finite-dimensional subspaces, usable stabilizer
information for pairs, and a limit procedure controlling both an
isometry and its inverse. None is proved here, and no transitive
non-Hilbertian separable example has been constructed.

\subsection{Sources}
\begin{itemize}
\item[S1]On Mazur rotations problem and its multidimensional versions.
\url{https://arxiv.org/abs/2012.08344}.
\textit{Formulation source.}
\end{itemize}
\end{document}
